% Part: first-order-logic
% Chapter: models-theories
% Section: set-theory

\documentclass[../../../include/open-logic-section]{subfiles}

\begin{document}

\olfileid{fol}{mat}{set}

\olsection{Verzamelingenleer}

Vrijwel de gehele wiskunde kan in de verzamelingenleer worden
ontwikkeld. Daarvoor zijn verschillende zaken nodig. Ten eerste is een
stelsel van axioma's voor de relatie~$\in$ nodig. Er zijn verscheidene
axiomasystemen ontwikkeld, die soms onderling strijdige eigenschappen
van~$\in$ hebben. Eén stelsel neemt een bijzondere plaats in:
$\Log{ZFC}$, de Zermelo--Fraenkel-verzamelingenleer met het keuzeaxioma.
Dit is verreweg het meest gebruikte en bestudeerde stelsel, omdat de
axioma's ervan vrijwel alles blijken te bewijzen wat wiskundigen willen
kunnen bewijzen. Voordat dit kan worden vastgesteld, moet echter eerst
duidelijk zijn hoe we alles wat wiskundigen willen zeggen überhaupt
kunnen \emph{uitdrukken}. De taal bevat om te beginnen geen
!!{constant}s of !!{function}s. Daardoor is op het eerste gezicht niet
duidelijk hoe we over bepaalde verzamelingen (zoals $\emptyset$ of
$\Nat$) of over bewerkingen op verzamelingen (zoals $X \cup Y$ en
$\Pow{X}$) kunnen spreken, laat staan over constructies die ook andere
zaken dan verzamelingen omvatten, zoals relaties en functies.

Om te beginnen is ``is een element van'' niet de enige relatie die ons
interesseert: ``is een deelverzameling van'' lijkt bijna even
belangrijk. We kunnen ``is een deelverzameling van'' echter
\emph{definiëren} met behulp van ``is een element van''. Daarvoor
moeten we in de taal van de verzamelingenleer !!a{formula}~$!A(x, y)$
vinden waaraan een paar verzamelingen~$\tuple{X, Y}$ dan en slechts dan
voldoet als $X \subseteq Y$. Nu is $X$ precies dan een deelverzameling
van $Y$ als alle !!{element}s van~$X$ ook !!{element}s van~$Y$ zijn.
We kunnen $\subseteq$ dus definiëren met de formule
\[
\lforall[z][(z \in x \lif z \in y)]
\]
Telkens wanneer we de relatie~$\subseteq$ in een formule willen
gebruiken, kunnen we haar vervolgens door deze formule vervangen
(waarbij $x$ en $y$ passend worden vervangen en de gebonden
variabele~$z$ zo nodig wordt hernoemd). De extensionaliteit van
verzamelingen houdt bijvoorbeeld in dat verzamelingen~$x$ en $y$ die
elk in de andere zijn bevat, dezelfde verzameling moeten zijn: $x$ en
$y$ moeten dan dezelfde verzameling zijn. Dit kan
worden uitgedrukt door $\lforall[x][\lforall[y][((x \subseteq y \land y
    \subseteq x) \lif x = y)]]$, of, als we $\subseteq$ door de
bovenstaande definitie vervangen, door
\[
\lforall[x][\lforall[y][((\lforall[z][(z \in x \lif z \in y)] \land
    \lforall[z][(z \in y \lif z \in x)]) \lif x = y)]].
\]
Dit is inderdaad een van de axioma's van $\Log{ZFC}$: het
``extensionaliteitsaxioma''.

Er is geen !!{constant} voor $\emptyset$, maar we kunnen ``$x$ is
leeg'' uitdrukken door $\lnot \lexists[y][y \in x]$. De bewering
``$\emptyset$ bestaat'' wordt dan de !!{sentence}~$\lexists[x][\lnot
\lexists[y][y \in x]]$. Dit is een ander axioma van~$\Log{ZFC}$.
(Merk op dat uit het extensionaliteitsaxioma volgt dat er maar één lege
verzameling is.) Wanneer we in de taal van de verzamelingenleer over
$\emptyset$ willen spreken, schrijven we dus ``er is een verzameling
die leeg is en \dots'' Om bijvoorbeeld uit te drukken dat $\emptyset$
een deelverzameling van iedere verzameling is, kunnen we schrijven
\[
\lexists[x][(\lnot\lexists[y][y \in x] \land \lforall[z][x \subseteq
    z])]
\]
waarbij $x \subseteq z$ uiteraard op zijn beurt door zijn definitie
moet worden vervangen.

Om over bewerkingen op verzamelingen te spreken, zoals $X \cup Y$ en
$\Pow{X}$, moeten we een vergelijkbare werkwijze volgen. De taal van de
verzamelingenleer bevat geen functiesymbolen, maar we kunnen de
functionele relaties $X \cup Y = Z$ en $\Pow{X} = Y$ uitdrukken door
\begin{align*}
& \lforall[u][((u \in x \lor u \in y) \liff u \in z)]\\
& \lforall[u][(u \subseteq x \liff u \in y )]
\end{align*}
want de !!{element}s van $X \cup Y$ zijn precies de verzamelingen die
!!{element}s van~$X$ of !!{element}s van~$Y$ zijn, en de !!{element}s
van $\Pow{X}$ zijn precies de deelverzamelingen van~$X$. Dit staat ons
echter niet toe $x \cup y$ of $\Pow{x}$ te gebruiken alsof het termen
zijn: we kunnen alleen de volledige !!{formula}s gebruiken die de
relaties $X \cup Y = Z$ en $\Pow{X} = Y$ definiëren. We weten zelfs nog
niet of ooit aan deze relaties wordt voldaan, dat wil zeggen of
verenigingen en machtsverzamelingen altijd bestaan. De !!{sentence}
$\lforall[x][\lexists[y][\Pow{x} = y]]$ is bijvoorbeeld een ander
axioma van~$\Log{ZFC}$ (het machtsverzamelingsaxioma).

Hoe kunnen we nu over geordende paren of functies spreken? Daarvoor
moeten we uitleggen hoe geordende paren en functies als bijzondere
soorten verzamelingen kunnen worden opgevat. Het geordende paar
$\tuple{x, y}$ kan bijvoorbeeld worden gedefinieerd als de verzameling
$\{\{x\}, \{x, y\}\}$. Net als hiervoor kunnen we echter geen
!!a{function} invoeren die deze verzameling benoemt; we kunnen alleen
de relatie $\tuple{x, y} = z$ definiëren, dat wil zeggen
$\{\{x\}, \{x, y\}\} = z$:
\[
\lforall[u][(u \in z \liff (\lforall[v][(v \in u \liff v = x)] \lor
  \lforall[v][(v \in u \liff (v = x \lor v = y))]))]
\]
Dit zegt dat de !!{element}s~$u$ van~$z$ precies de verzamelingen zijn
die ofwel alleen $x$ als !!{element} hebben, ofwel alleen $x$ en~$y$ als
!!{element}s hebben (met andere woorden: de verzamelingen die gelijk
zijn aan $\{x\}$ of aan $\{x, y\}$). Hiermee kunnen we vervolgens
andere zaken uitdrukken, bijvoorbeeld dat $X \times Y = Z$:
\[
\lforall[z][(z \in Z \liff \lexists[x][\lexists[y][(x \in
        X \land y \in Y \land \tuple{x, y} = z)]])]
\]

Een functie $f \colon X \to Y$ kan worden opgevat als de relatie $f(x)
= y$, dus als de verzameling paren~$\Setabs{\tuple{x,y}}{f(x) = y}$.
We kunnen dan zeggen dat een verzameling~$f$ een functie van $X$ naar
$Y$ is als (a) zij een relatie $\subseteq X \times Y$ is, (b) zij
totaal is, dat wil zeggen dat er voor iedere $x \in X$ een $y \in Y$
bestaat waarvoor $\tuple{x, y} \in f$, en (c) zij functioneel is, dat
wil zeggen dat telkens wanneer $\tuple{x, y}, \tuple{x, y'} \in f$,
$y = y'$ (want functiewaarden moeten uniek zijn). De bewering ``$f$ is
een functie van $X$ naar $Y$'' kan dus als volgt worden geschreven:
\begin{align*}
 \lforall[u][(u \in f \lif {}] & \lexists[x][\lexists[y][(x \in X \land y \in
      Y \land \tuple{x, y} = u)]]) \land {}\\
 \lforall[x][(x \in X \lif {}] &
      (\lexists[y][(y \in Y \land \mathrm{maps}(f, x, y))] \land {}\\
& (\lforall[y][\lforall[y'][((\mathrm{maps}(f, x, y) \land
    \mathrm{maps}(f, x, y')) \lif y = y')]]))
\end{align*}
waarbij $\mathrm{maps}(f, x, y)$ een afkorting is voor
$\lexists[v][(v \in f \land \tuple{x, y} = v)]$ (deze !!{formula}
drukt ``$f(x) = y$'' uit).

Nu is het evenmin moeilijk om bijvoorbeeld uit te drukken dat
$f\colon X \to Y$ !!{injective} is:
\begin{multline*}
 f \colon X \to Y \land \lforall[x][\lforall[x'][((x \in X \land x' \in
     X \land {}]] \\
 \lexists[y][(\mathrm{maps}(f, x, y) \land \mathrm{maps}(f,
        x', y))]) \lif x = x')
\end{multline*}
Een functie~$f\colon X \to Y$ is !!{injective} dan en slechts dan als,
wanneer $f$ de elementen $x, x' \in X$ op eenzelfde~$y$ afbeeldt,
$x = x'$. Als we deze formule afkorten als $\mathrm{inj}(f, X, Y)$,
kunnen we in de taal van de verzamelingenleer reeds een niet-triviale
uitspraak als de stelling van Cantor formuleren: er bestaat geen
!!{injective} functie van $\Pow{X}$ naar $X$:
\[
\lforall[X][\lforall[Y][(\Pow{X} = Y \lif
    \lnot\lexists[f][\mathrm{inj}(f, Y, X)])]]
\]

Men zou kunnen denken dat de verzamelingenleer nog een axioma nodig
heeft dat voor iedere bepalende eigenschap het bestaan van een
verzameling waarborgt. Als $!A(x)$ een formule van de
verzamelingenleer is waarin de variabele~$x$ vrij voorkomt, kunnen we
de !!{sentence}
\[
\lexists[y][\lforall[x][(x \in y \liff !A(x))]].
\]
beschouwen. Deze !!{sentence} zegt dat er een verzameling~$y$ bestaat
waarvan de !!{element}s precies de $x$ zijn die aan~$!A(x)$ voldoen.
Dit schema heet het ``comprehensieprincipe''. Het lijkt zeer bruikbaar,
maar is helaas inconsistent. Neem $!A(x) \ident \lnot x \in x$; het
comprehensieprincipe zegt dan
\[
\lexists[y][\lforall[x][(x \in y \liff x \notin x)]],
\]
dat wil zeggen: het beweert dat er een verzameling bestaat van alle
verzamelingen die geen !!{element}s van zichzelf zijn. Zo'n verzameling
kan niet bestaan---dit is de paradox van Russell. $\Log{ZFC}$ bevat
dan ook een beperkte---en consistente---versie van dit principe: het
separatieprincipe:
\[
\lforall[z][\lexists[y][\lforall[x][(x \in y \liff (x \in z \land
      !A(x))]]].
\]

\begin{prob}
Toon aan dat het comprehensieprincipe inconsistent is door
!!a{derivation} te geven die aantoont dat
\[
\lexists[y][\lforall[x][(x \in y \liff x \notin x)]] \Proves \lfalse.
\]
Het kan helpen om eerst $(A \lif \lnot A) \land (\lnot A \lif A)
\Proves \lfalse$ aan te tonen.
\end{prob}
\end{document}
